\medskip

On considère la fonction $f$ définie sur $]0~;~+\infty[$ par 
$f(x) = x^2 - 8\ln (x)$,
où ln désigne la fonction logarithme népérien.

On admet que $f$ est dérivable sur $]0~;~+\infty[$, on note $f'$ sa fonction dérivée. 

\medskip

\begin{enumerate}
\item On détermine $\displaystyle\lim_{x \to 0} f(x)$.

$\left .
\begin{array}{l}
\ds\lim_{x \to 0} x^2 =0\\
\ds\lim_{x \to 0 \atop x>0} \ln(x)=-\infty
\end{array}
\right \}
\text{ donc }
\ds\lim_{x \to 0\atop x>0} x^2 -8\ln(x)=+\infty
\text{ et donc }
\ds\lim_{x \to 0\atop x>0} f(x)=+\infty$


\item On admet que, pour tout $x > 0$,\: $f(x) = x^2\left(1 - 8\dfrac{\ln (x)}{x^2}\right)$.


$\left .
\begin{array}{l}
\ds\lim_{x \to +\infty} \dfrac{\ln(x)}{x} =0\\
\ds\lim_{x \to +\infty} \dfrac{1}{x}=0 
\end{array}
\right \}
\text{ donc }
\ds\lim_{x \to +\infty} \dfrac{\ln(x)}{x^2} =0
\text{ et donc }
\ds\lim_{x \to +\infty} \left ( 1-\dfrac{\ln(x)}{x^2}\right ) =1$

$\left .
\begin{array}{l}
\ds\lim_{x \to +\infty}  \left ( 1-\dfrac{\ln(x)}{x^2}\right ) =1\\
\ds\lim_{x \to +\infty} x^2=+\infty
\end{array}
\right \}
\text{ donc }
\ds\lim_{x \to +\infty}  x^2\left ( 1-\dfrac{\ln(x)}{x^2}\right ) =+\infty
\text{ et donc }
\ds\lim_{x \to +\infty}  f(x) =+\infty$

\item 
Calcul de la dérivée sur $]0~;~+\infty[$:

$f'(x)=2x -8\times \dfrac{1}{x}=\dfrac{2x^2-8}{x}=\dfrac{2\left (x^2-4\right )}{x}
=\dfrac{2(x-2)(x+2)}{x}$

\item Pour étudier les variations de $f$ sur $]0~;~+\infty[$, on cherche le signe de $f'(x)$.

\begin{center}
{\renewcommand{\arraystretch}{1.5}
\def\esp{\hspace*{3cm}}
$\begin{array}{|c |l *{4}{c} |} 
\hline
x  & 0 & \esp & 2 & \esp  & +\infty \\
\hline
x-2 &  & \pmb{-} &  \vline\hspace{-2.7pt}{0} & \pmb{+} &   \\
\hline
x+2 &  & \pmb{+} &  \vline\hspace{-2.7pt}{\phantom 0} & \pmb{+} &   \\
\hline
x &0  & \pmb{+} &  \vline\hspace{-2.7pt}{\phantom 0} & \pmb{+} &   \\
\hline
f'(x) &\vline\;\vline  & \pmb{-} &  \vline\hspace{-2.7pt}{0} & \pmb{+} &   \\
\hline
\end{array}$
}
\end{center}

$f(2)=4-8\ln(2)$
On en déduit le tableau des  variations de la fonction $f$  sur $]0~;~+\infty[$.

\begin{center}
{\renewcommand{\arraystretch}{1.3}
\psset{nodesep=3pt,arrowsize=2pt 3}
\def\esp{\hspace*{0.8cm}}
\def\hauteur{20pt}
$\begin{array}{|c|l*8{c}|}
\hline
x & 0  & \esp &  & \esp & 2 & \esp &  & \esp & +\infty \\ 
\hline
\text{ Signe de $f'$ } &\vline\;\vline  & &   \pmb{-}  & & \vline\hspace{-2.7pt}0   &&  \pmb{+} &&\\ 
\hline
 & \vline\;\vline \Rnode{max1}{~+\infty} & &&&&  &     &  & \Rnode{max2}{+\infty}   \\  
\text{ Variations de $f$ } & \vline\;\vline &  &   && & &   &  &  \rule{0pt}{\hauteur} \\ 
 & \vline\;\vline &   &  &  &  \Rnode{min}{4-8\ln(2)} & & & &  \rule{0pt}{\hauteur}    
\ncline{->}{max1}{min} \ncline{->}{min}{max2}
 \\ 
\hline
\end{array} $
}
\end{center}


\item 
$4-8\ln(2) \approx -1,55 <0$ donc on complète le tableau de variations.

\begin{center}
{\renewcommand{\arraystretch}{1.3}
\psset{nodesep=3pt,arrowsize=2pt 3}
\def\esp{\hspace*{0.8cm}}
\def\hauteur{20pt}
$\begin{array}{|c|l*8{c}|}
\hline
x & 0  & \esp & \Rnode{alpha}{\alpha} & \esp & 2 & \esp & & \esp & +\infty \\ 
\hline
 & \vline\;\vline \Rnode{max1}{~+\infty} & &&&&  &     &  & \Rnode{max2}{+\infty}   \\  
\text{ Variations de $f$ } & \vline\;\vline &  &  \Rnode{zero1}{0} && & &   &  &  \rule{0pt}{\hauteur} \\ 
 & \vline\;\vline &   &  &  &  \Rnode{min}{4-8\ln(2)} & & & &  \rule{0pt}{\hauteur}    
 \ncline{max1}{zero1}  \ncline{->}{zero1}{min}
  \ncline{min}{max2}
  \ncline[linecolor=red,linestyle=dashed]{zero1}{alpha}
 \\ 
\hline
\end{array} $
}
\end{center}

La fonction $f$ est dérivable sur $]0\;;\;+\infty[$ donc continue sur cet intervalle, donc continue sur $]0\;;\;2[$.
Sur l'intervalle $]0\;;\;2[$, la fonction $f$ est continue et strictement décroissante; elle passe d'une valeur positive à une valeur négative donc, d'après le corollaire du théorème des valeurs intermédiaires, l'équation $f(x)=0$ admet une solution unique dans $]0\;;\;2[$. 
On l'appelle $\alpha$.


\item On admet que, sur l'intervalle $[2~;~ +\infty[$, l'équation $f(x) = 0$ admet une solution unique $\beta$. On complète le tableau de variations de $f$ et on en déduit le signe de $f$ sur $]0\;;\;+\infty[$.

\begin{center}
{\renewcommand{\arraystretch}{1.3}
\psset{nodesep=3pt,arrowsize=2pt 3}
\def\esp{\hspace*{0.8cm}}
\def\hauteur{20pt}
$\begin{array}{|c|l*8{c}|}
\hline
x & 0  & \esp & \Rnode{alpha}{\alpha} & \esp & 2 & \esp & \Rnode{beta}{\beta} & \esp & +\infty \\ 
\hline
 & \vline\;\vline \Rnode{max1}{~+\infty} & &&&&  &     &  & \Rnode{max2}{+\infty}   \\  
\text{ Variations de $f$ } & \vline\;\vline &  &  \Rnode{zero1}{0} && & &\Rnode{zero2}{0}   &  &  \rule{0pt}{\hauteur} \\ 
 & \vline\;\vline &   &  &  &  \Rnode{min}{4-8\ln(2)} & & & &  \rule{0pt}{\hauteur}    
 \ncline{max1}{zero1}  \ncline{->}{zero1}{min}
  \ncline{min}{zero2}  \ncline{->}{zero2}{max2}
  \ncline[linecolor=red,linestyle=dashed]{zero1}{alpha}
    \ncline[linecolor=red,linestyle=dashed]{zero2}{beta}
 \\ 
\hline
\text{ Signe de $f$ } & \vline\;\vline & \pmb{+} & \vline\hspace{-2.7pt}0 & &  \pmb{-} &  &\vline\hspace{-2.7pt}0 &  \pmb{+} &\\ 
\hline
\end{array} $
}
\end{center}


\item Pour tout réel $k$, on considère la fonction $g_k$ définie sur $]0~;~+\infty[$ par:
$g_k(x) = x^2 - 8\ln (x) + k$; donc $g_k(x)=f(x)+k$.

La fonction $g_k$ a donc les mêmes variations que $f$ et elle a donc pour minimum sur $]0~,~+\infty[$ le nombre $4-8\ln(2)+k$.
Pour que $g_k$ soit positive ou nulle, il faut que ce minimum soit positif ou nul, donc $4-8\ln(2)+k\geqslant 0$ soit $k\geqslant 8\ln(2)-4$.

Donc $8\ln(2)-4$ est la plus petite valeur de $k$ telle que $g_k\geqslant 0$ sur $]0\;;\;+\infty[$.


\end{enumerate}

\bigskip
